\section{Central frustration and conditional concentration}
\label{sec:main-results}

Throughout this section, $d\geq2$, $m=2(d-1)$, and $e$ is the full-degree
link fixed in Section~\ref{sec:setup}.  Let $z\in\cZ(G)\setminus\{\one\}$ and
write
\begin{equation}\label{eq:L-bstar}
  L=\dist(\one,z),\qquad b_*=\frac{m-1}{m}L.
\end{equation}

\subsection{Exact realization by opposite links}

\begin{lemma}[Private opposite links]\label{lem:private-opposites}
For every plaquette $p$ incident to $e$, the edge of $p$ opposite and parallel
to $e$ belongs to no other plaquette incident to $e$.  These $m$ opposite
links are pairwise distinct.
\end{lemma}

\begin{proof}
An incident plaquette is determined by a transverse direction
$\nu\in\{2,\ldots,d\}$ and a sign $\sigma\in\{-1,+1\}$.  Its opposite edge
is parallel to direction $1$ and translated from $e$ by
$\sigma\mathbf e_\nu$.  Distinct pairs $(\nu,\sigma)$ give distinct
translations.  A translated direction-$1$ edge can share a unit square with
$e$ only in the translation direction and with the specified sign.  Hence it
belongs to exactly one member of the incident family.
\end{proof}

\begin{proposition}[Central staple surjectivity]\label{prop:central-surjectivity}
For any prescribed tuple $(z_p)_{p\ni e}\in\cZ(G)^m$, there is a
center-valued assignment of all links other than $e$ such that
$S_p=z_p$ for every $p\ni e$.
\end{proposition}

\begin{proof}
Set all nonselected links initially to $\one$.  For each incident plaquette,
change only its private opposite link from Lemma~\ref{lem:private-opposites}.
Choose $z_p$ or $z_p^{-1}$ according to the orientation with which that link
enters the staple.  Centrality makes cyclic ordering irrelevant, and the
private-link property prevents the change from altering another incident
staple.  The resulting staple tuple is the prescribed one.
\end{proof}

The proposition is an analytic incidence proof.  Appendix~\ref{app:incidence}
records an independent exact matrix certificate: after deleting the selected
edge column, the $m$ opposite-edge columns form an identity matrix over every
coefficient ring, and in particular over $\mathbb F_2$ for $G=\SUtwo$.

\subsection{The conditional kernel}

Choose a target plaquette $q\ni e$.  Apply
Proposition~\ref{prop:central-surjectivity} with
\begin{equation}\label{eq:one-frustrated-pattern}
 S_p=\one\quad(p\neq q),\qquad S_q=z^{-1}.
\end{equation}
The inverse is a convention chosen so that the target holonomy is
$z^{-1}U$; replacing $z$ by $z^{-1}$ gives the alternative orientation.

\begin{proposition}[One-link conditional law]\label{prop:conditional-law}
At the boundary assignment \eqref{eq:one-frustrated-pattern}, the continuous
conditional kernel of $U_e$ is
\begin{equation}\label{eq:conditional-law}
 \gamma_t^z(\dd U)=Z_t(z)^{-1}
 p_t(U)^{m-1}p_t(z^{-1}U)\,\Haar,
\qquad
 Z_t(z)=\int_Gp_t(U)^{m-1}p_t(z^{-1}U)\,\Haar.
\end{equation}
The target plaquette is good exactly on
\begin{equation}\label{eq:target-good}
  \mathcal G_\rho^z=\{U:\dist(z,U)<\rho\}.
\end{equation}
\end{proposition}

\begin{proof}
Every plaquette factor not incident to $e$ is constant under the
conditioning and cancels between numerator and normalization.  The incident
factors have the staple form \eqref{eq:staple-form}.  Substitution of
\eqref{eq:one-frustrated-pattern}, centrality, and inversion symmetry of
$p_t$ give the integrand in \eqref{eq:conditional-law}.  Bi-invariance gives
$\dist(\one,z^{-1}U)=\dist(z,U)$, proving \eqref{eq:target-good}.  A fully
oriented derivation appears in Appendix~\ref{app:conditional-kernel}.
\end{proof}

\subsection{Heat-kernel Laplace principle}

We use the following standard logarithmic form of the small-time heat-kernel
asymptotic.  The inclusion of the cut locus is essential here.

\begin{lemma}[Uniform Varadhan asymptotic]\label{lem:varadhan}
For the heat kernel of $\tfrac12\Delta$ on compact $G$,
\begin{equation}\label{eq:uniform-varadhan}
 \lim_{t\downarrow0}\sup_{U\in G}
 \left|2t\log p_t(U)+\dist(\one,U)^2\right|=0.
\end{equation}
Consequently, for fixed staples $s=(s_1,\ldots,s_m)$, the normalized measures
with density proportional to $\prod_{i=1}^m p_t(Us_i)$ satisfy a large
deviation principle at speed $(2t)^{-1}$ and rate
\begin{equation}\label{eq:generic-staple-energy}
 F_s(U)-\min_VF_s(V),qquad
 F_s(U)=\sum_{i=1}^m\dist(\one,Us_i)^2.
\end{equation}
The estimates are uniform when $s$ ranges over a compact set.
\end{lemma}

The logarithmic heat-kernel limit is due to Varadhan
\cite{Varadhan1967Fundamental,Varadhan1967Diffusion}.  The normalized product
form follows from a compact-space Laplace argument; details, including the
uniform parameter statement, are given in Appendix~\ref{app:laplace}.

For \eqref{eq:conditional-law}, the unnormalized energy is
\begin{equation}\label{eq:Fz}
  F_z(U)=(m-1)\dist(\one,U)^2+\dist(z,U)^2.
\end{equation}
Its minimizers can be described without coordinates.

\begin{lemma}[Two-point metric barycenter]\label{lem:metric-barycenter}
The minimum of \eqref{eq:Fz} is
\begin{equation}\label{eq:Fz-min}
  \min_UF_z(U)=\frac{m-1}{m}L^2.
\end{equation}
Its complete minimizing set is
\begin{equation}\label{eq:minimizer-set}
 \mathcal M_z=
 \left\{U:\dist(\one,U)=\frac Lm,
              \ \dist(z,U)=\frac{m-1}{m}L\right\}.
\end{equation}
Equivalently, $\mathcal M_z$ consists of the points at arclength $L/m$ on
all minimizing geodesics from $\one$ to $z$.
\end{lemma}

\begin{proof}
Put $a=\dist(\one,U)$ and $b=\dist(z,U)$.  The triangle inequality gives
$a+b\geq L$, while the exact identity
\begin{equation}\label{eq:weighted-square}
 (m-1)a^2+b^2
 =\frac{m-1}{m}(a+b)^2
  +\frac1m\bigl((m-1)a-b\bigr)^2
\end{equation}
gives the lower bound in \eqref{eq:Fz-min}.  Equality holds precisely when
$a+b=L$ and $b=(m-1)a$, hence at the distances in
\eqref{eq:minimizer-set}.  Compactness gives a minimizing geodesic from
$\one$ to $z$, and its point at arclength $L/m$ realizes equality.  Equality
in the triangle inequality characterizes membership in a minimizing
geodesic.
\end{proof}

\subsection{Main obstruction theorem}

\begin{theorem}[Central-frustration obstruction]\label{thm:central-obstruction}
Let $G$, $z$, $d$, $m$, and $L$ be as above.  On $B_d$ there is an admissible
center-valued assignment of all links other than $e$ for which the
conditional law is \eqref{eq:conditional-law}.  The family
$(\gamma_t^z)_{t>0}$ satisfies a large deviation principle with rate
\begin{equation}\label{eq:normalized-rate}
  I_z(U)=F_z(U)-\frac{m-1}{m}L^2,
\end{equation}
and concentrates on $\mathcal M_z$.

For every $0<\rho<b_*=(m-1)L/m$,
\begin{equation}\label{eq:exact-good-rate}
 \lim_{t\downarrow0}2t\log\gamma_t^z(\mathcal G_\rho^z)
 =-m(b_*-\rho)^2<0.
\end{equation}
In particular,
\begin{equation}\label{eq:bad-limit}
 \lim_{t\downarrow0}
 \sup_{\omega}
 \gamma_{e,t}^{\omega}
 \bigl(\mathcal G_{q,\rho}(\omega)^c\bigr)=1.
\end{equation}
Thus no estimate of the form \eqref{eq:uniform-quantifier} can hold.
\end{theorem}

\begin{proof}
The construction and the conditional density follow from
Propositions~\ref{prop:central-surjectivity} and
\ref{prop:conditional-law}.  Lemma~\ref{lem:varadhan} gives the large
deviation principle, and Lemma~\ref{lem:metric-barycenter} gives its
normalization and zero set.

It remains to compute the cost of the good event.  For $U$ with
$b=\dist(z,U)\leq\rho$, the triangle inequality gives
$\dist(\one,U)\geq L-b$, hence
\begin{align}
 F_z(U)
 &\geq (m-1)(L-b)^2+b^2 \notag\\
 &=\frac{m-1}{m}L^2+m(b-b_*)^2.\label{eq:good-lower-bound}
\end{align}
Because $\rho<b_*$, the right side decreases as $b$ increases through
$[0,\rho]$.  A minimizing geodesic from $\one$ to $z$ contains points with
$b$ arbitrarily close to $\rho$ from below, so
\begin{equation}\label{eq:good-infima}
 \inf_{\mathcal G_\rho^z}F_z
 =\inf_{\closure{\mathcal G_\rho^z}}F_z
 =\frac{m-1}{m}L^2+m(b_*-\rho)^2.
\end{equation}
The large-deviation lower bound on the open set and upper bound on its
closure coincide, proving \eqref{eq:exact-good-rate}.  Hence the good
probability tends to zero and the bad probability under the constructed
boundary tends to one.  The supremum in \eqref{eq:bad-limit} is at least this
value and at most one.
\end{proof}

\begin{remark}[What is and is not unique]\label{rem:nonunique-minimizers}
The distance pair in \eqref{eq:minimizer-set} is unique.  The group element
need not be.  In $\SUtwo$ with $z=-\one$, the endpoint is at the cut locus
and $\mathcal M_z$ is a conjugacy two-sphere.  No step of the proof selects a
geodesic or assumes an injectivity-radius condition.
\end{remark}
